%%%%%%%%%%  scan idx 302  |  folio 289  %%%%%%%%%%
\SGAPrintPage{289}%
Suppose now that we are in situation (2.11.1) (resp. (2.11.2)), with \(f\) of finite tor-dimension. Then (III 4.8.1 and 5.8 b)) the functor \(Rf_*\) sends \(D(X)_{\uparf}\) to \(D(Y)_{\uparf}\), and therefore defines a homomorphism
\sgaStableTarget{sga6.e04.eq.2-12-3}\begin{equation}\tag{2.12.3}
 f_*:K^{\bullet}(X)\longrightarrow K^{\bullet}(Y).
\end{equation}
This is not in general a ring homomorphism. However, for \(y\in K_{\bullet}(Y)\) (resp. \(K^{\bullet}(Y)\)) and \(x\in K^{\bullet}(X)\), we have the formula
\sgaStableTarget{sga6.e04.eq.2-12-4}\begin{equation}\tag{2.12.4}
 f_*(f^*(y)x)=y f_*(x),
\end{equation}
where \(f^*\) is the homomorphism (2.12.1) (resp. (2.7.1)). The proof is analogous to that of (2.11.1.2). If \(g:Y\to Z\) is a morphism satisfying the same hypotheses as \(f\), the homomorphisms (2.12.3) satisfy the transitivity formula:
\sgaStableTarget{sga6.e04.eq.2-12-3-1}\begin{equation}\tag{2.12.3.1}
 g_*f_*=(gf)_*.
\end{equation}

The homomorphism (2.12.3) is one of the essential ingredients of the Riemann--Roch theorem, in both Algebraic Geometry and Analytic Geometry (for indications in the analytic case, see (Exposé 0, 5)).

Suppose that \(X\) and \(Y\) are separated and divisorial schemes (II 2.2.5) (as is the case, for example, when \(X\) and \(Y\) possess ample invertible sheaves). Then (2.9 (i)) \(K^{\bullet}(X)\) (resp. \(K^{\bullet}(Y)\)) is identified with \(K^{\bullet}(X)_{\naif}\) (resp. \(K^{\bullet}(Y)_{\naif}\)), and (2.12.3) therefore defines a homomorphism
\sgaStableTarget{sga6.e04.eq.2-12-3.occ2}\begin{equation}\tag{2.12.3'}
 f_*:K^{\bullet}(X)_{\naif}\longrightarrow K^{\bullet}(Y)_{\naif}.
\end{equation}
As already explained in (Exposé 0, 3.1), the definition of (2.12.3') was by no means evident a priori. Indeed, from the standpoint of our Seminar, it is one of the main motivations for the sorites of Exposé I through Exposé IV.

%%%%%%%%%%  scan idx 303  |  folio 290  %%%%%%%%%%
\SGAPrintPage{290}%
\sgaStableTarget{sga6.e04.sec.3}\section*{3. Supplements on the Grothendieck groups of schemes\protect\SGAAuthorFootnote{1}{The results of this section will not be used in Exposés V through VIII.}}

\sgaStableTarget{sga6.e04.prop.3-1-0}\textbf{Proposition 3.1.0.} Let
\[
\begin{tikzcd}
X \arrow[d,"f"'] & X' \arrow[l,"g"']\arrow[d,"f'"]\\
Y & Y' \arrow[l,"h"']
\end{tikzcd}
\]
be a Cartesian square of schemes, where \(Y\) is quasi-compact and \(f\) quasi-compact and quasi-separated, and let \(E\in\Ob D^b(X)\) be a complex with quasi-coherent cohomology. Suppose \(X\) and \(Y'\) tor-independent over \(Y\) (III 1.5). Suppose moreover either \(h\) of finite tor-dimension or \(E\) of finite flat amplitude rel. to \(f\). Under these conditions, the base-change arrow (cf. SGA 4 XVII 4)
\[
 Lh^*\,Rf_*E\longrightarrow Rf'_*\,Lg^*E
\]
is defined and is an isomorphism.

\uline{Proof}: To define the arrow, we use the trivial duality formula (SGA 4 XVII 2.3) for \(h\), which says that there is a canonical functorial isomorphism
\begin{equation}\tag{$*$}
 \Hom(Lh^*L,M)\xrightarrow{\sim}\Hom(L,Rh_*M)
\end{equation}
for \(L\in\Ob D^-(Y)\), \(M\in\Ob D^+(Y')\). Using (III 3.5.2) and (III 3.7.1), we easily verify that \(Lg^*E\in\Ob D^b(X')\) and \(Rf_*E\in\Ob D^-(Y)\), so, applying (*) with \(L=Rf_*E\) and \(M=Rf'_*Lg^*E\), we are reduced to defining an arrow
\[
 Rf_*E\longrightarrow Rf_*Rg_*Lg^*E.
\]
For this we take the image under \(Rf_*\) of the adjunction arrow
\[
 E\longrightarrow Rg_*Lg^*E,
\]
which comes from the trivial duality formula for \(g\). The base-change arrow

%%%%%%%%%%  scan idx 304  |  folio 291  %%%%%%%%%%
\SGAPrintPage{291}%
is therefore defined; it remains to show that it is an isomorphism. For this, we may first suppose \(Y\) and \(Y'\) affine. Then, by the Leray spectral sequence for an open covering of \(X\), we reduce to the case where \(X\) and \(X'\) are affine. We therefore have a cocartesian square of rings
\[
\begin{tikzcd}
B\arrow[r] & B'\\
A\arrow[u]\arrow[r] & A'\arrow[u],
\end{tikzcd}
\]
and we may suppose that \(E\) is defined by a bounded-above complex of \(B\)-modules. The tor-independence of \(A'\) and \(B\) over \(A\), which is written
\[
 B\otimes_A^{\mathrm L}A'\xrightarrow{\sim}B',
\]
then implies that in \(D(A')\) we have
\[
 E\otimes_A^{\mathrm L}A'\xrightarrow{\sim}E\otimes_B^{\mathrm L}B',
\]
whence the proposition.

Translating (3.1.0) into the language of Grothendieck groups, we obtain:

\sgaStableTarget{sga6.e04.prop.3-1-1}\textbf{Proposition 3.1.1.} \uline{Let}
\[
\begin{tikzcd}
X \arrow[d,"f"'] & X' \arrow[l,"g"']\arrow[d,"f'"]\\
Y & Y' \arrow[l,"h"']
\end{tikzcd}
\]

%%%%%%%%%%  scan idx 305  |  folio 292  %%%%%%%%%%
\SGAPrintPage{292}%
\uline{be a Cartesian square of schemes}. \uline{Suppose} \(Y\) \uline{and} \(Y'\) \uline{quasi-compact}, \(X\) \uline{and} \(Y'\) \uline{tor-independent over} \(Y\) (III 1.5), \(f\) \uline{proper}, \uline{and pseudo-coherent}, \uline{and the finiteness conjecture} (III 2.1) \uline{satisfied for} \(f\) \uline{and} \(f'\) (\uline{for example} \(Y\) \uline{and} \(Y'\) \uline{noetherian}, \uline{or} \(f\) \uline{projective}). \uline{Then}, \uline{if} \(f\) \uline{is perfect} (\uline{resp}. \(h\) \uline{of finite tor-dimension}), \uline{we have}, \uline{for} \(x\in K^{\bullet}(X)\) (\uline{resp}. \(K_{\bullet}(X)\)),
\[
 h^* f_*(x)=f'_* g^*(x),
\]
\uline{where} \(h^*,g^*\) \uline{are defined by} (2.7.1) (\uline{resp}. (2.12.1)), \uline{and} \(f_*,f'_*\) \uline{by} (2.12.3) (\uline{resp}. (2.11.1.1)) (\uline{note that by} (III 4.7.1) \(f'\) \uline{is perfect if} \(f\) \uline{is}, \uline{and that by} (III 3.5.1) \(g\) \uline{has finite tor-dimension if} \(h\) \uline{does}).

\sgaStableTarget{sga6.e04.para.3-2}3.2. \uline{Passage to the limit for} \(K^{\bullet}_{\naif}\) \uline{and} \(K_{\bullet}\) Let \((X_i,f_{ij})_{i\in I}\) be a filtered projective system of quasi-compact and quasi-separated schemes, with affine transition morphisms (cf. EGA IV 8.2.2). Put \(X=\varprojlim X_i\), and denote by \(f_i:X\to X_i\) the canonical homomorphism.

\sgaStableTarget{sga6.e04.para.3-2-1}3.2.1. The homomorphisms (2.7.2) \(f_i^*:K^{\bullet}(X_i)_{\naif}\longrightarrow K^{\bullet}(X)_{\naif}\) define, by the transitivity formula (2.7.3), a homomorphism
\sgaStableTarget{sga6.e04.eq.3-2-1-1}\begin{equation}\tag{3.2.1.1}
 \varinjlim K^{\bullet}(X_i)_{\naif}\longrightarrow K^{\bullet}(X)_{\naif}.
\end{equation}
I claim that \uline{the homomorphism} (3.2.1.1) \uline{is an isomorphism}. Indeed, let \(E\) be a locally free \(\underline O_X\)-Module of finite type. By (EGA IV 8.5.2 and 8.5.5), there is an index \(i\) and a locally free \(\underline O_{X_i}\)-Module of finite type \(E_i\) such that \(E\simeq f_i^*(E_i)\). By (EGA IV 8.5.2.5), the image of \(\cl(E_i)\) in \(\varinjlim K^{\bullet}(X_i)_{\naif}\) is independent of the choice of \(E_i\); denote it by \(u(E)\). By (EGA IV 8.5.2 and 8.5.6), the map
\(u:\Ob\Loclib(X)\longrightarrow \varinjlim K^{\bullet}(X_i)_{\naif}\) thus defined is additive (1.5);

%%%%%%%%%%  scan idx 306  |  folio 293  %%%%%%%%%%
\SGAPrintPage{293}%
it therefore defines a homomorphism \(K^{\bullet}(X)_{\naif}\longrightarrow \varinjlim K^{\bullet}(X_i)_{\naif}\), which is immediately seen, by (EGA IV 8.5.2), to be inverse to (3.2.1.1), proving our assertion.

\sgaStableTarget{sga6.e04.para.3-2-2}3.2.2. In the same way, one defines a canonical homomorphism
\sgaStableTarget{sga6.e04.eq.3-2-2-1}\begin{equation}\tag{3.2.2.1}
 \varinjlim K^{\bullet}(X_i)\longrightarrow K^{\bullet}(X).
\end{equation}
Using Deligne's localization technique (SGA 4 XVII), which implies that the category of perfect complexes on \(X\) is equivalent to the category \(\mathop{\underrightarrow{\mathrm{Lim}}}\) of the analogous categories \(\Parf(X_i)\), one again finds that (3.2.2.1) is an isomorphism.

\sgaStableTarget{sga6.e04.para.3-2-3}3.2.3. Suppose the \(X_i\) and \(X\) \uline{noetherian}, and the morphisms \(f_{ij}\) \uline{flat} (the \(f_i\) then are as well). Then the homomorphisms (2.12.1)
\(f_i^*:K_{\bullet}(X_i)\longrightarrow K_{\bullet}(X)\) define an isomorphism
\begin{equation}\tag{3.2.3.1} \varinjlim K_{\bullet}(X_i)\xrightarrow{\sim}K_{\bullet}(X).
\end{equation}
The proof is analogous to that of (3.2.1), using the fact (2.4) that \(K_{\bullet}(X)\simeq K_{\bullet}(X)_{\naif}\) is computed in terms of coherent sheaves on \(X\).

Since the module structures of the \(K_{\bullet}(X_i)\) over the \(K^{\bullet}(X_i)_{\naif}\) are compatible with the inverse images \(f_{ij}^*\) by (2.7 a)) and (2.12.2), we find, by passing to the limit, that \(\varinjlim K_{\bullet}(X_i)\) has a module structure over \(\varinjlim K^{\bullet}(X_i)_{\naif}\); this is compatible, via the isomorphisms (3.2.1.1) and (3.2.3.1) with the module structure of \(K_{\bullet}(X)\) over \(K^{\bullet}(X)_{\naif}\). In other words, we have a commutative diagram, where the vertical arrows are defined by (3.2.1.1) and (3.2.3.1):
\[
\begin{tikzcd}[column sep=tiny]
\varinjlim K^{\bullet}(X_i)_{\naif}\otimes \varinjlim K_{\bullet}(X_i)
\arrow[r]\arrow[d,"\simeq"]
&\varinjlim K_{\bullet}(X_i)\arrow[d,"\simeq"]\\
K^{\bullet}(X)_{\naif}\otimes K_{\bullet}(X)\arrow[r]
&K_{\bullet}(X).
\end{tikzcd}
\]

